\SourceLeaf{384}%
war viel geringer, als sie bei einer als regelmäßig vorausgesetzten Erdgestalt sein sollte. Die neue Messung hat gezeigt, dass \textsc{Beccaria} bei der astronomischen Bestimmung dieser Krümmung tatsächlich einen Fehler von \(13''41\) begangen hat (der bei der Unvollkommenheit seiner Instrumente sehr verzeihlich ist). Das Vorzeichen dieses Fehlers ist aber das Gegenteil des vermuteten, und die Anomalie vergrößert sich somit noch um diesen Betrag. Die aus den geodätischen Messungen mit den oben angegebenen Elementen berechnete Amplitude beträgt nämlich (auf \textsc{Beccarias} Endpunkte bezogen) \(1^0\,8'\,18''91\), die aus \textsc{Beccarias} astronomischen Beobachtungen folgende \(1^0\,7'\,44''30\) und die neu bestimmte \(1^0\,7'\,31''07\). Die neue Messung wurde mit so guten Hilfsmitteln und mit so außerordentlicher Sorgfalt durchgeführt, dass man sich gezwungen sieht, diesen großen Unterschied von \(47''84\) fast ganz als eine Unregelmäßigkeit der Erdgestalt anzusehen, die bemerkenswerteste Tatsache dieser Art, die bisher in den Annalen der höheren Geodäsie vorgekommen ist. Höchstwahrscheinlich ist die Anziehung der Alpenketten, die diese Messung im Norden und Süden begrenzen, eine Hauptursache dieses Phänomens. Ebenso wahrscheinlich hat die ungleiche Dichte der tieferen Erdschichten, vielleicht bis in große Tiefe hinab, einen mindestens ebenso großen Anteil daran. Zumindest lassen sich ähnliche, sehr beträchtliche Unterschiede zwischen ganz in der Ebene liegenden Punkten (z.\,B. eine Anomalie von \(21''9\) zwischen Mailand und Parma) kaum anders erklären. Wir fügen hinzu, dass sich mit der zunehmenden Zahl sorgfältig ausgeführter Gradmessungen immer mehr die Überzeugung durchsetzt, dass solche Abweichungen nur ihrer Größe nach als Ausnahmen gelten dürfen und dass sie sich in größerem oder kleinerem Maßstab überall zeigen. Die Verfasser haben eine interessante vergleichende Übersicht über die astronomisch bestimmten und die aus geodätischen Messungen berechneten Polhöhen an 34 Punkten gegeben, die über halb Europa verstreut und durch Dreiecke miteinander verbunden sind. Freilich kann man sie nur als einen unvollkommenen Versuch ansehen, da sie zum größten Teil lediglich auf unbeglaubigten, bruchstückhaften Mitteilungen über die Ergebnisse der geodätischen Messungen beruht. Leider sind nämlich die meisten dieser Messungen in Frankreich, Italien, Österreich und Bayern immer noch nicht veröffentlicht.

Derselbe Abschnitt enthält außerdem eine neue Messung einer kurzen Basislinie bei Turin, die einige Umstände der von Herrn~%
